% Source: 15 - Wed Oct 28 2026.pdf, page 4. Content remains in source order.
\sourcepage{15}{4}
PROOF: Let $y=y_i$ and let $\mathrm{GC}_i$ the greedy choice at iteration $i$. Clearly, $\mathrm{GC}_i=y_j$ for some $j<i\Rightarrow y_j=\mathrm{GC}_i\leq y=y_i$. Also, since $y=y_i\notin L'$ it must be $\mathrm{GC}_i\geq y_i/(1+\delta)$. The proof follows setting $z=\mathrm{GC}_i\in L'$.

EXERCISE: Prove that $L'=\mathrm{TRIM}(L,\delta)$ is a minimum trimming (standard analysis of a greedy algorithm).
However, we will not use minimality directly but the following:
\begin{notesresult}
\textbf{PROPERTY} Let $L'=\langle z_0=0,z_1,\ldots,z_{K-1}\rangle=\mathrm{TRIM}(L,\delta)$.
Then, for $1\leq j<K-1$: $z_{j+1}/z_j>(1+\delta)$.
(The property does not hold for $j=0$ since $z_0=0$.)
\end{notesresult}
PROOF. $z_{j+1}$ and $z_j$ are two consecutive greedy choices. Thus it must be
\[z_j<z_{j+1}/(1+\delta)\Rightarrow z_{j+1}/z_j>(1+\delta).\]
\subsection{The approximation algorithm}
We are ready to transform EXP-SS$(\langle S,t\rangle)$ into a FPTAS:
(identical to EXP-SS$(S,t)$ but for trimming with $\delta=\eps/(2n)$).
\begin{notesalgorithm}
\textbf{APPROX-SS ($S,t,\eps$)}
\begin{notespseudocode}
\pseudocodeline{0}{* Let $S=\{x_1,x_2,\ldots,x_n\}$ *}
\pseudocodeline{0}{$L'_0\gets\langle\varnothing\rangle$}
\pseudocodeline{0}{\textbf{for} $i\gets1$ \textbf{to} $n$ \textbf{do}}
\pseudocodeline{1}{$L'_i\gets\mathrm{MERGE\text{-}WD}(L'_{i-1},L'_{i-1}+_t x_i)$}
\pseudocodeline{1}{$L'_i\gets\mathrm{TRIM}(L'_i,\eps/(2n))$}
\pseudocodeline{0}{\textbf{return} $\mathrm{MAX}(L'_n)$}
\end{notespseudocode}
\end{notesalgorithm}
